\scanpage{276}
\mspage{226}
mais surtout ils permettent de définir, pour \unsure{$S$|$C$} de $M$
\[
\left\{\begin{aligned}
W_S^g&\overset{\text{déf}}{=}W_S^{\mathrm{Liss}}\subset\mathrm{Fl}(M_{/S})\\
W_S^d&=W_S^{\mathrm{Prop}}\subset\mathrm{Fl}(M_{/S})
\end{aligned}\right.
\]
avec les propriétés suivantes :

\deleted{\textcircled{\unsure{b}}}\textcircled{a} Si $f:S'\to S$ est dans $\mathrm{Prop}$ \inserted{(resp. $\mathrm{Liss}$)}, alors
\[
\begin{aligned}
f^*:M_{/S}\longrightarrow M_{/S'}\quad &\text{transforme }W_S^d\text{ en }W_{\unsure{S|S'}}^d\\
&(\text{resp. } W_S^g\text{ en }W_{\unsure{S|S'}}^g)
\end{aligned}
\]
ce qui est une pure tautologie, et ne concerne pas encore les relations de $\mathrm{Prop}$ avec $\mathrm{Liss}$, mais par contre :
\[
\text{\textcircled{b}}\quad
\left\{\begin{aligned}
(\underline{\mathrm{Liss}}\,S)(W_S^{\unsure{u|v}})^{-1}&\xrightarrow{\approx}(M_{/S})(W_S^d)^{-1}
\quad\text{une équivalence}\\
\text{et}\qquad W_S^u&=W_S^d\quad\text{sur }\underline{\mathrm{Liss}}\,S
\end{aligned}\right.
\]
et dualement
\[
\left\{\begin{aligned}
(\underline{\mathrm{Prop}}\,S)(W_S^u)^{-1}&\xrightarrow{\approx}(M_{/S})(W_S^g)^{-1}
\quad\text{une \unsure{équivce|équiv\textsuperscript{ce}}}\\
\text{et}\qquad W_S^u&=W_S^g\quad\text{sur }\underline{\mathrm{Prop}}\,S.
\end{aligned}\right.
\]
De plus, il y a lieu d'énoncer \inserted{\textcircled{c}} \deleted{\unsure{une|une fois}} \inserted{deux} propriétés d'adjonction \inserted{(duales l'une de l'autre)}, pour lesquelles je renvoie à XII, théorème \unsure{salade} 2, p. 108-113. Il se pose ici la question subsidiaire de simplifier les conditions de validité de la formule d'adjonction, un peu hétéroclites dans loc. cit.

Il se pose \deleted{\unsure{en troi}} d'autre part la question de faire la théorie des \deleted{la} lois covariantes $f_!$ et $f_{\unsure{*|\star}}$ pour $\mathrm{Ho}^{\unsure{\mathrm{lc}|\mathrm{li}}}_W(S)$ de $=$ pour $\mathrm{Ho}^{\mathrm{Prop}}_W(S)$, $\mathrm{Ho}^{\mathrm{Liss}}_W(S)$.
% archive sheet number: 276
